\documentclass[11pt]{article}
\usepackage{fontspec}
\setmainfont{Times New Roman}
\setsansfont{Arial}
\setmonofont{Consolas}
\usepackage{amsmath,amssymb,mathtools}
\usepackage{geometry}
\usepackage{microtype}
\usepackage{enumitem}
\geometry{a4paper,margin=2.28cm}
\setlength{\parindent}{1.15em}
\setlength{\parskip}{0pt}
\makeatletter
\renewcommand\footnotesize{\@setfontsize\footnotesize{7.5}{8.5}\selectfont}
\makeatother
\emergencystretch=120em
\allowdisplaybreaks
\sloppy
\hbadness=10000
\vbadness=10000
\hfuzz=2pt
\vfuzz=10pt
\raggedbottom

\begin{document}

% Unit: Paper33 complete v001. Direct Interslavic Latin translation from refreshed Zenodo R122 Paper32/Paper33 source slice, Paper33 lines 525--569 / segments P33-S0001--P33-S0010. The older local Paper33 final-audited slice is superseded for this unit because it omits the R122/source-scan opening of the two-problem paragraph.
\section*{33. Hiperkompleksne veličiny i teorija predstavjenj v aritmetičnom poimanju}
\emph{Atti Congresso Bologna 2 (1928), s. 71--73.}

\begin{center}
\textsc{Emmy Noether} (Göttingen -- Germania)\\[-0.1em]
\rule{3em}{0.4pt}\\[0.8em]
\textsc{Hiperkompleksne veličiny i teorija predstavjenj}\\
\textsc{v aritmetičnom poimanju}\footnote{Podrobno izloženje bude javiti se v \emph{Math. Zeitschr.}; v prodolženju tam bude razsmotrjena takože nekomutativna Galoisova teorija. Sr. takože o “razpadnyh poljah”: R. Brauer i E. Noether, \emph{Ber. Berl. Ak.} 1927, a o občih strukturnyh teoremah -- doklad G. Köthe v tih \emph{Atti}.}
\end{center}

Hoču pokazati, kako se teorija predstavjenj -- osobito predstavjenja grup -- može razuměti kako teorija klasov modulov i idealov; i kako se te poslědnje znovu podčinjejut obobčenomu pojmu grupy, imenno grupam s operatorami.\footnote{Grupy s operatorami izhode od W. Krull (\emph{Math. Zeitschr.} 23) i O. Schmidt (\emph{Math. Zeitschr.} 29). I tut dano obrabotanje problema redukcije danogo predstavjenja izhodi -- v ograničenju na komutativne těla predstavjenja -- od W. Krull (\emph{Theorie und Anwendung der verallgemeinerten Abelschen Gruppen}, Heidelberger Ber. 1926).}

Pod \emph{predstavjenjem grupy} -- v danom polju $P$ -- razuměje se, kako vědomo, slědujuče: elementam $a,b,\ldots$ grupy $G$ pridružajut se matrice $A,B,\ldots$ s koeficientami iz $P$ tako, že produktu $ab$ odgovarja produkt $AB$; sistem matric stavaje se homomorfnym obrazom grupy. S tym predstavjenjem dano jest jednovrěmenno predstavjenje “grupnogo kolca”, ktoro sostoji iz vsih “grupnyh čisel”, t. j. vsih linearnyh kombinacij $a\alpha+b\beta+\cdots+c\gamma$, kde $a,b,\ldots,c$ proběgajut vse kombinacije iz konečno mnogih grupnyh elementov, a $\alpha,\beta,\ldots,\gamma$ sootvětno iz konečno mnogih elementov iz $P$; pri tom $a,b,\ldots,c$ smatrajut se linearně nezavisnymi, a množenje jest dano grupnym množenjem i računskymi zakonami. Predstavjenje grupnogo kolca jest jednovrěmenno homomorfno vzhodno k složenju; tako predstavjenja grup podčinjajut se občemu problemu \emph{predstavjenja kolc}, kde potrěbna jest homomorfnost vzhodno k složenju i množenju.

Tut sut suštna dva problema. Prvy jest problem \emph{redukcije danogo} predstavjenja i pytanje o jednoznačnosti ireducibilnyh sostavnyh česti, ktore pri tom nastavajut. Drugi sostoji v pytanju o cělosti možlivyh predstavjenj -- može byti samo ireducibilnyh -- danogo kolca v danom, ne nužno komutativnom tělu.

Prvy problem jest mnogo prostši; to se vidi uže iz togo, že ni o kolcu, ktoro treba predstaviti, ni o -- v obče nekomutativnom -- tělu predstavjenja ne sut potrěbne nikake specialne predpostavjenja. Fakt, že pri predstavjenju ide o matrice fiksovanogo stepena, daje vse potrěbne predpostavjenja konečnosti. Problem možno popolno obrabotati posrědstvom teorije grup s operatorami, na ktoru zato kratko prějdu.

Grupa $\mathfrak{M}$ nazyvaje se \emph{grupa s oblastju operatorov}, ako jednovrěmenno dany jest sistem simbolov $\Theta,H,\ldots$ -- operatorov -- tako, že kombinacije $\Theta(a),H(a),\ldots$ za vsako $a$ iz $\mathfrak{M}$ sut jednoznačno opreděljene, dajut element iz $\mathfrak{M}$ i zadovoljajut distributivny zakon: $\Theta(ab)=\Theta(a)\Theta(b)$. Dvě grupy s jednakoju oblastju operatorov nazyvajut se operatorno-izomorfne, ako sut izomorfne v obyčnom smyslu i ako, krome togo, iz odgovarjanja $a$ i $\bar a$ slěduje takože odgovarjanje $\Theta(a)$ i $\Theta(\bar a)$, $H(a)$ i $H(\bar a)$, itd. Ako potom kako podgrupy dopuskajut se samo take, ktore same dopuskajut oblast operatorov, tako že grupa nazyvaje se prosta, ako ne imaje dopustimogo vlastnogo normalnogo dělitelja kromě jedinicy, togda Jordan-Hölderov teorem o kompozicijnom redu ostavaje važny v formě: ako grupa s oblastju operatorov vobče imaje kompozicijny red, togda sistem faktor-grup jest v smyslu operatornogo izomorfizma jednoznačno opreděljen do porjadka. Pod toj samoju predpostavkoju važi takože teorem o razkladu, jednoznačnom v smyslu operatornogo izomorfizma, na direktno nerazložime faktory.

Vez s problemom predstavjenj dana jest tym, že k grupam s operatorami prinadležet osobito idealy i moduly, smatrjane kako abelove grupy vzhodno k složenju, poky operatori sut dany množenjem s elementami kolca. Vsako predstavjenje poraždaje se modulom predstavjenja, t. j. modulom vzhodno k kolcu i tělu predstavjenja. Točnije, vsaky klas modulov, t. j. klas medžu soboju operatorno-izomorfnyh modulov predstavjenja, poraždaje to samo predstavjenje. Tako pytanje o jednoznačnosti redukcije predstavjenja razrěšaje se teoremom o kompozicijnom redu i teoremom o direktnom produktu. Kompozicijny red, priměnjeny k modulu predstavjenja, daje za vsaku matricu redukciju formy
\[
\begin{pmatrix}
B_1&0&0\\
&B_2&0\\
A_{ik}&\cdots&B_r
\end{pmatrix},
\]
pri čem diagonalne matrice, sootvětny faktor-grupam, poraždajut “ireducibilno” predstavjenje, t. j. tako, ktoro samo uže ne dopuskaje taku redukciju.

No poneže klasy tih faktor-grup sut jednoznačno opreděljene, to samo važi za diagonalne predstavjenja, jednoznačno v smyslu ekvivalencije predstavjenj. Analogično teorem o direktnom produktu daje jednoznačnost “razpada” na “nerazložime” sostavne česti:
\[
\begin{pmatrix}
C_1&&&\\
&C_2&&\\
&&\cdots&\\
&&&C_s
\end{pmatrix},
\]
pri čem vsako $C_i$ može potom ješče byti jednoznačno reducirano po teoremu o kompozicijnom redu.

\emph{Drugi problem} hoču načrtati v slučaju predstavjenja hiperkompleksnyh sistemov, obče -- kolc, ktore zadovoljajut dvojno uslovje cěpov. Rezultat jest, že vsaky ireducibilny (prosty) klas modulov jest klas idealov; zato kompozicijne redy iz jednostrannyh idealov črez svoje faktor-grupy uže daju vse ireducibilne predstavjenja. Točnije, dostatočne sut kompozicijne redy v kolcu klasov ostatkov po radikalu, ktoro stavaje se “popolno reducibilnym” kolcom. Problem jest tym svedeny k strukturnomu izslědovanju popolno reducibilnyh kolc; i zato pokazuje se sootvětnym problemu brati predstavjenja v tělahi avtomorfizmov kako izhodišče. To treba razuměti tako: vsi prosti prave idealy dvustranno proste komponente takogo kolca prinadležet k tomu samomu klasu, zato imajut to samo kolco avtomorfizmov, ktoro zbog prostosti idealov stavaje se tělom, jednovrěmenno kolcom avtomorfizmov prostyh lěvyh idealov. Samo prosto kolco stavaje se izomorfno sistemu vsih matric nad tělom avtomorfizmov; i iz togo predstavjenja izhode vse predstavjenja vzhodno k podtělu, osobito vzhodno k oblasti koeficientov hiperkompleksnogo sistema, kogda same elementy těla avtomorfizmov zaměnjajut se pridruženymi matricami. Tako se tut Wedderburnov strukturny teorem, vědomy v slučaju hiperkompleksnyh sistemov, novo tolkuje črez pojem těla avtomorfizmov, ktory izhodi iz občej teorije grup, i jednovrěmenno poznaje se kako izvor teorije predstavjenj hiperkompleksnyh sistemov. Nastavajut nove pytanja Galoisovej teorije nekomutativnyh těl, ktore tesno svęzane sut s pytanjem o “razpadnyh poljah”. Jednovrěmenno pokazuje se vyhled k strukturnym teoremam občih, ne popolno reducibilnyh kolc posrědstvom kolc avtomorfizmov nerazložimyh idealov. Teorija grup, v svojem razširjenju na najobčejše grupy s operatorami, i tut znovu pokazuje se kako uredivajuči princip.

\end{document}
